\section{Relative complements of families}

\import{set/finite.tex}

For a family of subsets of a fixed ambient set, taking relative complements
turns unions into empty intersections.  This elementary set-theoretic bridge
is recorded here for later use in the open-cover characterization of
compactness.

\begin{definition}\label{relative_complement_family}
    $\relativeComplements{S}{F}=\{S\setminus A\mid A\in F\}$.
\end{definition}

\begin{proposition}\label{relative_complement_family_finite}
    Suppose $F$ is finite.
    Then $\relativeComplements{S}{F}$ is finite.
\end{proposition}
\begin{proof}
    Let $C=\{G\in\pow{F}\mid
        \text{$\relativeComplements{S}{G}$ is finite}\}$.
    $\relativeComplements{S}{\emptyset}=\emptyset$
        by \auto{relative_complement_family,emptyset,setext}.
    Hence $\emptyset\in C$
        by \obv{powerset_bottom,emptyset_is_finite}.
    Show that for every $G,A$ such that $G\in C$ and $A\in F$ we have
        $\cons{A}{G}\in C$.
    \begin{subproof}
        Fix $G,A$.
        Assume $G\in C$ and $A\in F$.
        $G\subseteq F$ by \obv{pow_iff,subseteq}.
        Hence $\cons{A}{G}\subseteq F$
            by \obv{cons_subseteq_intro}.
        Thus $\cons{A}{G}\in\pow{F}$ by \obv{pow_iff,subseteq}.
        $\relativeComplements{S}{\cons{A}{G}}
            =\cons{S\setminus A}{\relativeComplements{S}{G}}$
            by \auto{relative_complement_family,cons_iff,setext}.
        $\relativeComplements{S}{G}$ is finite.
        Hence $\relativeComplements{S}{\cons{A}{G}}$ is finite
            by \auto{finite_cons}.
    \end{subproof}
    Hence $F\in C$ by \obv{finite_induction}.
\end{proof}

\begin{proposition}\label{relative_complement_family_involutive}
    Suppose $F\subseteq\pow{S}$.
    Then $\relativeComplements{S}{\relativeComplements{S}{F}}=F$.
\end{proposition}
\begin{proof}
    Show that $\relativeComplements{S}{\relativeComplements{S}{F}}
        \subseteq F$.
    \begin{subproof}
        It suffices to show that every element of
            $\relativeComplements{S}{\relativeComplements{S}{F}}$
            is an element of $F$ by \obv{subseteq}.
        Fix $B\in\relativeComplements{S}{\relativeComplements{S}{F}}$.
        Take $D\in\relativeComplements{S}{F}$ such that
            $B=S\setminus D$ by \obv{relative_complement_family}.
        Take $A\in F$ such that $D=S\setminus A$
            by \obv{relative_complement_family}.
        $A\subseteq S$ by \obv{subseteq,pow_iff}.
        Hence $S\setminus(S\setminus A)=A$
            by \obv{double_relative_complement}.
        Thus $B=A$.
    \end{subproof}
    Show that $F\subseteq
        \relativeComplements{S}{\relativeComplements{S}{F}}$.
    \begin{subproof}
        It suffices to show that every element of $F$ is an element of
            $\relativeComplements{S}{\relativeComplements{S}{F}}$
            by \obv{subseteq}.
        Fix $A\in F$.
        $A\subseteq S$ by \obv{subseteq,pow_iff}.
        $S\setminus A\in\relativeComplements{S}{F}$
            by \obv{relative_complement_family}.
        $S\setminus(S\setminus A)=A$
            by \obv{double_relative_complement}.
        Follows by \obv{relative_complement_family}.
    \end{subproof}
    Follows by \obv{subseteq_antisymmetric}.
\end{proof}

\begin{proposition}\label{relative_complement_union_iff_intersection_empty}
    Suppose $F$ is inhabited and $F\subseteq\pow{S}$.
    Then $\unions{\relativeComplements{S}{F}}=S$ iff
        $\inters{F}=\emptyset$.
\end{proposition}
\begin{proof}
    $\unions{\relativeComplements{S}{F}}\subseteq S$
        by \auto{relative_complement_family,setminus_subseteq,unions_iff,subseteq}.
    $\inters{F}\subseteq S$
        by \auto{inters_iff_forall,subseteq,pow_iff}.
    Show that if $\unions{\relativeComplements{S}{F}}=S$, then
        $\inters{F}=\emptyset$.
    \begin{subproof}
        Assume $\unions{\relativeComplements{S}{F}}=S$.
        Show that $\inters{F}\subseteq\emptyset$.
        \begin{subproof}
            It suffices to show that every element of $\inters{F}$
                is an element of $\emptyset$ by \obv{subseteq}.
            Fix $x\in\inters{F}$.
            $x\in S$ by \obv{subseteq}.
            Hence $x\in\unions{\relativeComplements{S}{F}}$.
            Take $D\in\relativeComplements{S}{F}$ such that $x\in D$
                by \obv{unions_iff}.
            Take $A\in F$ such that $D=S\setminus A$
                by \obv{relative_complement_family}.
            $x\in A$ by \obv{inters_destr}.
            $x\notin A$ by \auto{setminus}.
            Contradiction.
        \end{subproof}
        Follows by \obv{subseteq_emptyset_iff}.
    \end{subproof}
    Show that if $\inters{F}=\emptyset$, then
        $S\subseteq\unions{\relativeComplements{S}{F}}$.
    \begin{subproof}
        Assume $\inters{F}=\emptyset$.
        It suffices to show that every element of $S$ is an element of
            $\unions{\relativeComplements{S}{F}}$ by \obv{subseteq}.
        Fix $x\in S$.
        Suppose not.
        Show that for every $A\in F$ we have $x\in A$.
        \begin{subproof}
            Fix $A\in F$.
            Suppose not.
            $x\in S\setminus A$ by \obv{setminus}.
            $S\setminus A\in\relativeComplements{S}{F}$
                by \obv{relative_complement_family}.
            Hence $x\in\unions{\relativeComplements{S}{F}}$
                by \auto{unions_iff}.
            Contradiction.
        \end{subproof}
        Hence $x\in\inters{F}$ by \auto{inters_iff_forall}.
        Thus $x\in\emptyset$.
        Contradiction by \obv{emptyset}.
    \end{subproof}
    Follows by \auto{subseteq_antisymmetric}.
\end{proof}
